\section{Introduction: Paper 17's wall and where the two spines meet}\label{sec:intro}

\subsection{The question \texorpdfstring{$(\T)$}{(T)} and the wall of Paper 17}\label{sec:wall-context}

Let $d \ge 2$, let $b = \clog{d}$ be the canonical entry bit-width, and let
$L = d^2 b = \Theta(d^2\log d)$ be the bit-length of a $d \times d$ matrix of $b$-bit words
laid out row by row on the read-only input tape of a deterministic multi-tape Turing machine
with $k$ read-write work tapes and a write-only output tape
(\cite[Definition~2.1]{FirstPaper2026}). Question $(\T)$---and its stronger time-space trade-off
form, the Tape Coding Conjecture \cite[Conjecture~9.1]{FirstPaper2026}---asks whether every
such machine that outputs the transposed matrix $X^\top$ for all inputs $X$ requires
$\omega(L)$ steps; by Harvey and van der Hoeven \cite{HvdH2025}, a negative answer to the
$O(L)$-time question implies $\Omega(n\log n)$ integer multiplication on multi-tape Turing
machines.

After sixteen papers of reductions and barriers, the synthesis paper
\emph{Order, Symmetry and the Swap} \cite[Section~4]{SeventeenthPaper2026} paused to diagnose why
every prior attack had stalled short of $(\T)$:
\begin{enumerate}[label=(\roman*)]
\item \emph{Motion without data} (\cite[Sections~4.2 and~4.3]{SeventeenthPaper2026}): single-scale
  pass counting, trace crossing bounds, and non-abelian commutator invariants saturate at one
  dyadic scale or $O(1)$ passes once a machine is permitted to interleave scales or code across
  head visits; even oblivious $O(L)$-time transposers need macroscopic crossings
  (\cite[Theorem~5.5, Corollary~5.6, and Section~10.3]{SixteenthPaper2026}), which a two-tape
  cascade can supply.
\item \emph{Codes without machines} (\cite[Section~4.4]{SeventeenthPaper2026}): once a computation
  is compressed to an abstract constant-depth or bounded-degree storage DAG whose nodes may apply
  \emph{arbitrary, distinct} functions at every tape cell, network coding defeats the
  $\frac{1}{2}$-rate pairing barrier. Paired-distance separation already exceeds rate $\frac{1}{2}$
  (\cite[Theorem~7.1, Theorem~7.3, and Corollary~8.2]{StorageCodes2026}, reaching
  $(\sqrt{3}-1)^2 - \eps \approx 0.536$ over $\F_{p^2}$) and reaches binary rate
  $1 - O(\Delta^{-\alpha}) \to 1$ as the degree cap $\Delta \to \infty$ across arbitrary odd-girth
  separation thresholds (\cite[Definition~5.1, Theorem~5.5, and Theorem~6.1]{HighRate2026},
  formalized conditionally on an odd-girth linear seed).
\item \emph{The missing fourth object} (\cite[Sections~4.1 and~4.6]{SeventeenthPaper2026}): a real
  Turing machine has a \emph{fixed finite transition table} $\delta$. When a head sweeps across
  regularly laid-out data, successive cells are not decoded by $d^2$ bespoke linear combinations
  tailored to each matrix coordinate $(i,j)$; they are decoded by the \emph{same} finite-state
  rule applied to shifted views of the tape.
\end{enumerate}

To turn that observation into theorems, \cite[Sections~5 and~6]{SeventeenthPaper2026} isolated two
algebraic pillars of shared rules on a vector space $\Mat_{d\times d}(F)$:
\emph{Order} (\cite[Theorem~5.1 and Corollary~5.5]{SeventeenthPaper2026}, showing that a two-term
swap identity $A_1 X B_1 + A_2 X B_2 = X^\top$ forces $\dim S \le d(d+1)/2$ via triangularization
of a matrix pencil over the algebraic closure) and \emph{Symmetry} (\cite[Definition~6.3 and
Theorem~6.4]{SeventeenthPaper2026}, showing that an $M$-class linear shared rule is governed by
cyclic column shifts $\rot_c(X)_{i,j} = X_{i,\,j+c}$ and satisfies the translate-pigeonhole
bound).
Two papers then pushed the two halves of the program forward independently:
\cite{AdaptiveTransfer2026} established macroscopic original-entry capacity and bounded-fan-in cut
pullback for adaptive machines, while \cite{SharedDecoders2026} developed the capacity and
multi-stage composition theory of $M$-class shared decoders.

\subsection{Two spines, two complementary blind spots}\label{sec:blind-spots}

Stepping back across \cite{SeventeenthPaper2026,AdaptiveTransfer2026,SharedDecoders2026} reveals
why neither spine could close $(\T)$ alone: each spine preserved one half of what makes a tape
sweep special and discarded the other.

\paragraph{The Machine Spine: bounded arity $\Delta = O(1)$ without decoder sharing.}
Trace compression via $\TCP(k)$ (\cite[Definition~2.3, Proposition~2.4, and
Theorem~7.5]{BoundedStack2026}) shows that cutting a $k$-head walk DAG with an $\eps$-fraction vertex
cut reduces residual path length to $D = D(k,\eps) = O_{k,\eps}(1)$, trapping each clean output cell
inside a causal cone of bounded read fan-in $\Delta \le b(k+1)^D$ (or $\Delta \le (k+1)^D$ per bit
or word-step; conditional on a depth-$D$ source cut $B$, \cite[Lemma~H.2, Theorem~H.5, and
Corollary~H.6]{AdaptiveTransfer2026} transports interior-lane cones across adaptive runs). Yet the
moment one forgets that those cones are evaluated by a single transition table $\delta$, two
barriers block an unconditional lower bound: (a)~abstract separated storage codes reach rate
$1 - O(\Delta^{-\alpha}) \to 1$ as $\Delta \to \infty$ (\cite{StorageCodes2026,HighRate2026}), and
(b)~any $O(n\log n)$-gate Boolean barrel-shift circuit embeds into a max-degree-$4$ causal network
achieving complete-exchange shift coding with $\Omega(1)$ coding-to-routing gain when vertices apply
distinct local gate functions (\cite[Propositions~11.2 and~11.3]{FirstPaper2026}).

\paragraph{The Shared-Decoder Spine: $M$-class sharing without bounded arity $\Delta = O(1)$.}
Conversely, \cite{SeventeenthPaper2026,SharedDecoders2026} enforced $M$-class decoder sharing:
column $c$ of $X^\top$ is produced by applying one of $M$ shared functions $f_{p(c)}$ to the
column-shifted state $\rot_c(X)$. However, in the core shared-decoder predicates of
\cite{SharedDecoders2026}, each $f_m$ was allowed \emph{unrestricted} read arity ($\Delta = d^2$).
Dropping $\Delta = O(1)$ erected three artificial walls:
\begin{enumerate}[label=\textbf{Wall \arabic*.},leftmargin=*]
\item \textbf{Arbitrary-set blow-up at unrestricted arity.}
  Fixing an arbitrary nonlinear subset $C \subseteq \Sigma^{d\times d}$ arises naturally when one
  freezes a trace cut (partitioning inputs into nonlinear trace/cut fibres). At unrestricted read
  arity $\Delta = d^2$, a single shared deterministic decoder ($M = 1$) already achieves rate
  $1 - \frac{\log_{\abs{\Sigma}} d}{d^2} \to 1$ via cyclic-orbit markers
  (\cite[Theorem~A.1 and Corollary~A.2]{SharedDecoders2026}). In \cite[Section~6.7,
  Propositions~6.21--6.23, and Problem~9.4(b)]{SeventeenthPaper2026}, written (unformalized)
  acyclic-subgraph and order arguments settled bounded offsets (\cite[Proposition~6.21(a)]{SeventeenthPaper2026}),
  window one ($\Delta = 1$, \cite[Proposition~6.21(b)]{SeventeenthPaper2026}), exact $1$-class
  window two ($M = 1, \Delta \le 2, \mathcal{E} = \emptyset$, rate $\le \frac{1}{2} + \frac{3}{2d}$,
  \cite[Proposition~6.22(b)]{SeventeenthPaper2026}), and bounded-value offsets or mutual helpers
  (\cite[Proposition~6.22(c)--(d) and Proposition~6.23]{SeventeenthPaper2026}). Yet for every
  constant window $\Delta \ge 3$ with spread rows and offsets (\cite[Problem~9.4(b)]{SeventeenthPaper2026}),
  as well as for sublinearly growing class counts $M = o(d/\Delta)$ or nonzero broken-output sets
  $\mathcal{E} \ne \emptyset$---especially over the canonical word alphabet
  $\abs{\Sigma_d} = 2^{\clog{d}} \ge d$, where one symbol can encode its own shift index $c$---no
  nontrivial upper bound $< 1$ on arbitrary sets was known.
\item \textbf{Hales--Jewett collapse in two-stage linear composition.}
  A two-pass machine induces a composition $E_2 \circ E_1 = \tau$ of two shared stages.
  In \cite[Section~6, Lemma~6.5, and Theorem~6.6]{SharedDecoders2026}, bounding the two-stage rank
  relied on extracting a monochromatic length-$r$ arithmetic progression in the stage-1 partition
  via the Hales--Jewett theorem \cite{HalesJewett1963}. Because Hales--Jewett numbers grow faster
  than primitive-recursively, that bound yields only $d - o(d)$ rank and collapses once $M_1$
  grows as slowly as $\log\log d$; worse, Behrend colorings \cite{Behrend1946} avoid $3$-term
  progressions with $M_1 = 2^{O(\sqrt{\log d})}$ classes, and random colorings cap monochromatic
  progressions at $O(\log d)$ even for $M_1 = 2$.
\item \textbf{Aperiodic and Stone-shuffle blindness.}
  As highlighted in \cite[Problems~9.1 and~9.5]{SeventeenthPaper2026}, the cyclic shift $\rot_c$
  couples only pairs $(i,j)$ and $(i',j')$ on the same anti-diagonal $i + j = i' + j'$ (when rows
  and shifts vary together across class-uniform collisions). Moreover,
  many shared-decoder invariants treat matrix transpose $\tau(X)_{i,j} = X_{j,i}$ and Stone's
  perfect shuffle/unshuffle permutation \cite{Stone1971} on an equal footing, even though Stone's
  shuffle and unshuffle are computable in $O(L)$ steps on multi-tape Turing machines ($3$ work tapes
  for shuffle and $2$ work tapes for unshuffle, without hardcoded $d$ or $b$) in $O(1)$ passes!
\end{enumerate}

\subsection{Overview of results and how they reshape the attack on \texorpdfstring{$(\T)$}{(T)}}\label{sec:results-overview}

This paper shows that reuniting \emph{bounded read arity} ($\Delta = O(1)$) with \emph{shared
decoders} ($M$ classes), replacing arithmetic progressions by \emph{affine Hilbert cubes} and
\emph{shifted sumsets}, and tracking \emph{Operator Schmidt Rank} across the row/column tensor
bipartition breaks through all three walls simultaneously, answers the concrete algebraic
questions of \cite[Remark~5.6 and Problem~9.3]{SeventeenthPaper2026}, and builds the native
transition-table bridge of \cite[Problem~9.4]{SeventeenthPaper2026}. \Cref{tab:overview}
summarizes the correspondence between the walls and open problems of
\cite{SeventeenthPaper2026,SharedDecoders2026} and our kernel-verified theorems.

\begin{table}[t]
\centering
\small
\renewcommand{\arraystretch}{1.22}
\begin{tabular}{@{}p{0.27\textwidth}p{0.43\textwidth}p{0.25\textwidth}@{}}
\toprule
\textbf{Prior barrier / problem} & \textbf{Kernel-verified resolution in this paper} & \textbf{Theorem / Lean module} \\
\midrule
Wall 1 \& \cite[Problem~9.4(b)]{SeventeenthPaper2026}: arbitrary sets reach rate $1 - o(1)$ at $\Delta = d^2$; spread windows $\Delta \ge 3$, $M = o(d/\Delta)$, and $\mathcal{E} \ne \emptyset$ open at rate $1$. &
Alphabet-independent capacity $\log_{\abs{\Sigma}}\abs{C} \le d^2 - (d^2 - \abs{\mathcal{E}} - dM\Delta)\frac{\Delta^\Delta}{(\Delta+1)^{\Delta+1}}$; Lean-verified rate $< 1$ for all $\Delta \ge 1$ and strict deficit for $M\Delta = o(d)$. &
\Cref{thm:set-capacity,cor:arity-two,thm:growing-deficit} (\leanfile{BoundedAritySetCapacityProofs.lean}) \\
\addlinespace[2pt]
Wall 2: two-stage linear Hales--Jewett collapse at $M_1 \approx \log\log d$ (\cite[Section~6]{SharedDecoders2026}). &
\emph{Theorem 2A} (any field $F$): affine $k$-cubes give exact-factorization cutoff $d \le \max(8M_1 M_2^2, (2M_1)^{2^{M_2}})$ ($M_1 = o(d^{1/2^{M_2}})$ for fixed $M_2$, covering Behrend sets). \emph{Theorem 2B} ($(d:F)\ne 0$): sumset annihilating polynomials + KST give $\Omega(bd(\log d)^{1/(2M_2-1)})$ rank deficit ($\Omega(bd\log d)$ at $M_2=1$). &
\Cref{thm:2a,thm:2b,cor:2b-m2-one} (\leanfile{AffineCubeCompositionProofs.lean}, \leanfile{SumsetRankDeficitProofs.lean}) \\
\addlinespace[2pt]
Wall 3 \& \cite[Problems~9.1, 9.5]{SeventeenthPaper2026}: $\rot$-invariance \& Stone-shuffle blindness. &
Anti-diagonal tensor untwisting $F^{d\times d} \cong F^d \otimes_F F^d$ + $5/8$ generic mask floor; Operator Schmidt Rank $\OSR(\tau)=d^2$ vs.\ $\OSR(S)=4$, forcing $t \ge \log_2 d$ passes on shuffle/row/column pipelines. &
\Cref{thm:anti-diagonal,thm:osr-pass,cor:general-osr-pass} (\leanfile{OperatorDualityProofs.lean}) \\
\addlinespace[2pt]
\cite[Remark~5.6]{SeventeenthPaper2026}: does $\sum_{\ell=1}^3 A_\ell X B_\ell = X^\top$ obey $\dim S \le d(d+1)/2$? &
\textbf{No} over every field: Kronecker family $K_{3m}$ achieves $4m(m+1)$ (exceeding $\frac{3m(3m+1)}{2}$ for $m \in \{1,2,3,4\}$), plus $11 > 10$ at $d = 4$ and $29 > 28$ at $d = 7$ ($d \in \{3,4,6,7,9,12\}$). &
\Cref{thm:three-term-kronecker,thm:three-term-unified} (\leanfile{SwapCodeAnswersProofs.lean}) \\
\addlinespace[2pt]
\cite[Problem~9.3]{SeventeenthPaper2026}: multi-scale dyadic masks on both sides of $X$. &
Interval bound $\frac{d^2 + d\min(\beta_A,\beta_B)}{2}$ with $\beta_\bullet = \max(1, 2\sum_j M_{\bullet,j})$ ($\frac{1}{2} + o(1)$ for sublinear profiles); Walsh boundary dichotomy ($2^{j+1}$ boundaries vs.\ $\le 2$). &
\Cref{thm:multiscale-interval,thm:walsh-dichotomy} (\leanfile{SwapCodeAnswersProofs.lean}) \\
\addlinespace[2pt]
\cite[Problem~9.4]{SeventeenthPaper2026}: native $\TM$ transition-table sweep extraction. &
Symbolic \lean{CausalConeTree} \& \lean{SlotConeTree} extraction from $M.\delta$ on frozen trace cuts + conditional reductions of $(\T)$. &
\Cref{thm:sweep-cut-absorbed,thm:sweep-cut-refined,thm:conditional-tcc} (\leanfile{MachineSweepConsequences.lean}) \\
\bottomrule
\end{tabular}
\caption{Summary of how the main results address the walls and open problems of
\cite{SeventeenthPaper2026,SharedDecoders2026}.}\label{tab:overview}
\end{table}

Let us highlight six aspects of these results and their exact scope:
\begin{enumerate}[label=\arabic*.]
\item \textbf{Bounded-arity arbitrary-set capacity (\Cref{sec:capacity}).}
  \Cref{thm:set-capacity} proves that for \emph{every} finite alphabet $\abs{\Sigma} \ge 2$,
  \emph{every} subset $C \subseteq \Sigma^{d\times d}$, and \emph{every} $M$-class deterministic
  shared decoder of cell read arity $\Delta \ge 1$ accurate outside $\mathcal{E} \subseteq
  \Z_d \times \Z_d$,
  \[
    \log_{\abs{\Sigma}}\abs{C} \;\le\; d^2 \;-\;
    \bigl(d^2 - \abs{\mathcal{E}} - dM\Delta\bigr)\,p_{\min}(\Delta),
    \qquad
    p_{\min}(\Delta) \;:=\; \frac{\Delta^\Delta}{(\Delta+1)^{\Delta+1}} = \Theta(1/\Delta).
  \]
  At $\Delta = 1$ with $\mathcal{E} = \emptyset$, \cite[Proposition~6.21(b)]{SeventeenthPaper2026}
  gives the written cycle-counting rate $\frac{1}{2} + \frac{M}{2d}$, and at $M = 1, \Delta \le 2,
  \mathcal{E} = \emptyset$, \cite[Proposition~6.22(b)]{SeventeenthPaper2026} gives a written
  row-ranking bound of rate $\frac{1}{2} + \frac{3}{2d}$. For every constant arity $\Delta \ge 3$ on
  spread relations (\cite[Problem~9.4(b)]{SeventeenthPaper2026}), as well as across sublinearly
  growing class counts $M = o(d/\Delta)$ and broken-output sets $\mathcal{E} \ne \emptyset$ at every
  $\Delta \ge 1$, \Cref{thm:set-capacity} is the \emph{first} nontrivial ($< 1$) capacity bound on
  arbitrary nonlinear sets---and the first Lean-verified bound $< 1$ at $\Delta \ge 2$, giving e.g.\
  $\le \frac{23}{27} + \frac{8}{27d}$ at $M=1, \Delta=2$ (\Cref{cor:arity-two}) and
  $\le \frac{229}{256} + o(1)$ at $\Delta = 3$. Crucially, the formula
  has \emph{zero dependence on $\abs{\Sigma}$}: even over the canonical word alphabet
  $\abs{\Sigma_d} = 2^{\clog{d}} \ge d$, where reading a single symbol can reveal the shift index $c$,
  a fixed read footprint $Q_{p(c),r}$ of size $\Delta < d$ misses the self-referential coordinate
  $(c, r-c)$ on at least $d - M\Delta$ shifts $c$, forcing the $1 - \Theta(1/\Delta)$ ceiling
  without counting functions over a fixed alphabet. Whether the $M = 1, \Delta = O(1)$ ceiling for
  $\Delta \ge 3$ drops all the way to $\frac{1}{2} + o(1)$ (the question of
  \cite[Problem~9.4(b)]{SeventeenthPaper2026}) or admits nonlinear coding gain above $\frac{1}{2}$
  remains open and is formalized as an exact logical dichotomy (\Cref{prop:half-barrier-dichotomy}).

\item \textbf{Affine Hilbert cubes and Ramsey-free two-stage cutoffs (\Cref{sec:cubes}).}
  Over an \emph{arbitrary} field $F$, we replace arithmetic progressions by affine $k$-cubes
  $c_0 + \{0,\delta_1\} + \cdots + \{0,\delta_k\} \subseteq \Z_d$ with distinct nonzero steps
  $\delta_r$. \Cref{lem:affine-cube} proves that an affine $k$-cube in one stage-1 class forces
  $(k+1)b$ independent linear constraints per stage-2 class on the common message $W$ (not a naive
  $2^k$-multiplier recurrence, which fails already at $k = 2$: at $d = 13, q = 2$,
  \Cref{ex:d13-sharpness} builds an explicit $3$-AP-free $2$-cube $J = \{0,1,4,5\} \subseteq \Z_{13}$
  attaining $\dim W = 9 > 8 = \lfloor(26 + 6)/4\rfloor$ over every field with $(13 : F) \ne 0$,
  while obeying $3\dim W \le 30 \implies \dim W \le 10$ over every field). Because affine $k$-cubes
  are bipartite box patterns extractable by iterated Cauchy--Schwarz (\Cref{lem:cs-cube}),
  \Cref{thm:2a} eliminates Hales--Jewett numbers and proves the explicit cutoff
  $d \le \max\bigl(8M_1 M_2^2,\; (2M_1)^{2^{M_2}}\bigr)$, ruling out exact two-stage linear
  factorizations for $M_1$ up to $\Theta(d^{1/2^{M_2}})$ at fixed $M_2$ (including $3$-AP-free
  Behrend partitions with $M_1 = 2^{O(\sqrt{\log d})}$) and $M_2$ up to $(1-\eps)\log_2\log_2 d$ at
  fixed $M_1$. As an essential coded-lane control at the boundary $d = 4 = M_1 M_2$,
  \Cref{ex:d4-coded} shows that even when $K_1$ and $K_2$ reuse the \emph{exact same} two-interval
  partition $p_1 = p_2 = [0,0,1,1]$ on $\Z_4$, exact full transposition $K_2 \circ K_1 = \tau$ holds
  for all lane widths $b \ge 1$ over every field $F$ with $(5 : F) \ne 0$ (including scalar $b = 1$
  over $\F_2$!), whereas uncoded entry-copy and bit-lane-copy factorizations are impossible for all
  $b \ge 1$ and $\abs{\Sigma} \ge 2$.

\item \textbf{Shifted sumsets and polynomial-in-$\log d$ two-stage rank deficits (\Cref{sec:sumsets}).}
  Over any field $F$ with $(d : F) \ne 0$, we go qualitatively beyond common-message growth for
  \emph{fixed} class counts $M_1, M_2$. On any shifted $(2M_2)$-fold monochromatic sumset
  $\widetilde{J}_0 + \widetilde{B}_1 + \cdots + \widetilde{B}_{2M_2-1}$ in one stage-1 class with
  $\abs{\widetilde{J}_0} = k$, $0 \in \widetilde{B}_t$, and $\abs{\widetilde{B}_t} = k + 1$, an
  injective arbitrary-subset Vandermonde map over roots of unity (\Cref{lem:vandermonde-subset})
  combines with a stacked band pullback (\Cref{lem:band-pullback}) and $2M_2 - 1$ sparse
  annihilating polynomials of degree $\le \diam_{\N}(\widetilde{B}_t)$ (\Cref{lem:sumset-poly}) to
  force $2\rank(\tau - K_2 K_1) \ge bk\bigl(d - \sum_t \diam_{\N}(\widetilde{B}_t)\bigr)$
  (\Cref{thm:sumset-rank}). Combining this with iterated K\H{o}v\'ari--S\'os--Tur\'an sumset
  extraction (\Cref{lem:kst-sumset}) proves \Cref{thm:2b}: for all constants $M_1, M_2 \ge 1$,
  every exact or approximate two-stage linear composition over $(d : F) \ne 0$ has error rank
  $\rank(\tau - K_2 K_1) = \Omega\!\bigl(b\,d\,(\log d)^{1/(2M_2-1)}\bigr)$, and $\Omega(b\,d\log d)$
  when $M_2 = 1$ (\Cref{cor:2b-m2-one}), compared to the sub-iterated-logarithmic $o(d)$ bound of
  \cite[Theorem~6.6]{SharedDecoders2026}. \Cref{ex:d12-band} proves over $\C$ at $d = 12$ why true
  stage-1 shift invariance cannot be relaxed to an abstract $1$-class band subspace.

\item \textbf{Anti-diagonal untwisting and Operator Schmidt Rank duality (\Cref{sec:duality}).}
  Conjugating by the anti-diagonal transform $Z_{c,a} = X_{c,\,a-c}$ intertwines $\rot_t$ with the
  pure right-factor shift $I_d \otimes \sigma^t$ on $F^d \otimes_F F^d$, decoupling
  shift-collision constraints into $d$ independent anti-diagonal slices (\Cref{thm:anti-diagonal}),
  packaging the $9/16$ and $5/8$ random-mask certificate floors of
  \cite[Appendix~C]{SharedDecoders2026}, and realizing the single-shift $5/8$ floor via generic
  tensor subspaces (\Cref{prop:mask-floors}). Across the row/column tensor cut
  $\End_F(F^d \otimes F^d) \cong \End_F(F^d) \otimes \End_F(F^d)$, Operator Schmidt Rank cleanly
  separates matrix transpose ($\OSR(\tau) = d^2$, \Cref{thm:osr-swap}) from Stone's perfect shuffle
  ($\OSR(S) = 4$ for even $d = 2^m$, \Cref{thm:osr-shuffle}), giving a $t \ge \log_2 d$ pass lower
  bound on pipelines of row-local, column-local, and Stone-shuffle layers (\Cref{thm:osr-pass})
  while certifying $O(L)$-time \leanname{Transpose.TM} implementations of $3$-tape shuffle, $2$-tape
  unshuffle, and $2$-tape row shear (\Cref{prop:stone-tm}). Two controls
  (\Cref{prop:shear-control,prop:otp-control}) delineate the exact scope of $\OSR$ on row-diagonal
  gauges and nonlinear one-time-pad encodings.

\item \textbf{Answers to Remark~5.6 and Problem~9.3 of Paper 17 (\Cref{sec:swapcodes}).}
  We resolve \cite[Remark~5.6]{SeventeenthPaper2026} in the negative over \emph{every} field $F$:
  three-term swap subspaces satisfying $A_1 X B_1 + A_2 X B_2 + A_3 X B_3 = X^\top$ attain exact
  dimension $4m(m+1)$ at $d = 3m$ for all $m \ge 1$ (strictly exceeding $\frac{3m(3m+1)}{2}$ for
  $m \in \{1,2,3,4\}$, i.e., $d \in \{3,6,9,12\}$), plus $11 > 10$ at $d = 4$ and $29 > 28$ at
  $d = 7$, exceeding $d(d+1)/2$ across $d \in \{3,4,6,7,9,12\}$ (\Cref{thm:three-term-unified}),
  with $d = 5$ ($15$ vs.\ $15$) isolated as a crisp open threshold. We also resolve the literal
  dyadic-interval formulation of \cite[Problem~9.3]{SeventeenthPaper2026}, proving the upper bound
  $(d^2 + d\min(\beta_A,\beta_B))/2$ with $\beta_\bullet = \max(1, 2\sum_j M_{\bullet,j})$ and full
  $d^2$ attainment at the complete dyadic profile (\Cref{thm:multiscale-interval}), and establish
  the structural boundary dichotomy separating contiguous dyadic intervals ($\le 2$ cyclic
  boundaries) from periodic Rademacher/Walsh masks ($2^{j+1}$ cyclic boundaries, singleton-atom
  intersection, yet $\le (m+1)2^m$ single-mask $\rot$-span; \Cref{thm:walsh-dichotomy}).

\item \textbf{Native transition-table sweep extraction on frozen cuts (\Cref{sec:sweep}).}
  Addressing \cite[Problem~9.4]{SeventeenthPaper2026}, we prove that on any frozen trace/cut fibre
  of a native \leanname{Transpose.TM} ($1 \le M.k$, extendable from $M.k = 0$ via dummy-tape
  padding \Cref{lem:pad-dummy-tape}) whose cut-free causal paths have length $\le D$, the machine's
  transition function $M.\delta$ evaluates explicit depth-$D$ symbolic \lean{CausalConeTree}s of
  cell support $\le b(M.k+1)^D$. Periodic relative tree templates induce a genuine
  \lean{IsBoundedAritySetCode} either by absorbing cut-touching outputs into $\mathcal{E}$
  (\Cref{thm:sweep-cut-absorbed}) or---with \emph{zero} extra broken cells---by refining $M_0$ via
  $L_{\mathrm{cut}}$ frozen cut-slot labels per column into
  $M_{\mathrm{cls}} = M_0 \bigl((M.q+1)(M.g+1)^{M.k}\bigr)^{L_{\mathrm{cut}}}$ classes
  (\Cref{thm:sweep-cut-refined}). We conclude in \Cref{sec:frontier} with conditional theorems
  showing precisely which remaining trace-extraction, half-barrier, and macroscopic rank hypotheses
  would close $(\T)$.
\end{enumerate}

\subsection{Formalization status and conventions}\label{sec:conventions}

\emph{Every} numbered theorem, lemma, proposition, corollary, and example in this paper is
completely formalized and kernel-verified in Lean~4 (version \texttt{v4.34.1} with Mathlib
\cite{Mathlib2026}) inside the \texttt{BoundedArityDecoders} library, depending solely on Lean's
three standard foundational axioms (\texttt{propext}, \texttt{Classical.choice}, and
\texttt{Quot.sound}), with zero \texttt{sorry}, \texttt{admit}, \texttt{native\_decide}, or custom
axioms. Each numbered result carries a \textsf{Lean:} header naming its formal declarations and
source module, followed by a self-contained mathematical proof outline faithful to the formal proof.
Results quoted in prose from earlier papers of the series record their exact manuscript locator
(noting explicitly whenever an upstream statement is written-only or conditional); where we
formalize the separated storage-code amplification benchmark from \cite{HighRate2026} in
\Cref{prop:diag-high-rate}(a), its formal statement retains the explicit seed hypothesis
\lean{HasOddGirthLinearSeed} (matching \leanname{TapeStorageRates.OddGirthLinearSeed} of
\cite[Definition~5.1]{HighRate2026}). Question $(\T)$ itself
remains open; all implications toward $(\T)$ in \Cref{sec:sweep} are stated with their exact formal
hypotheses.
